\documentclass[11pt,a4paper]{article}
\usepackage[T1]{fontenc}
\usepackage{lmodern,amsmath,amssymb,amsthm,mathtools,booktabs,microtype}
\usepackage[margin=25mm]{geometry}
\usepackage[colorlinks=true,linkcolor=blue,citecolor=blue,urlcolor=blue]{hyperref}
\setlength{\emergencystretch}{3em}
\newtheorem{theorem}{Theorem}[section]
\newtheorem{lemma}[theorem]{Lemma}
\newtheorem{proposition}[theorem]{Proposition}
\newtheorem{corollary}[theorem]{Corollary}
\newcommand{\Q}{\mathbb Q}\newcommand{\Z}{\mathbb Z}
\DeclareMathOperator{\N}{N}
\title{Smaller norm polynomials with a possibly singular final quadratic extension}
\author{}\date{15 November 2026}
\begin{document}\maketitle
\begin{abstract}
A separating element in a quadratic tower remains separating on geometric
characters when the final radicand is allowed to vanish. Galois descent
then recovers a rational character from a rational determinant root even
on a nonreduced fiber. The pole valuation obstruction permits a zero final
ordinate, so this removes two large norm factors from each of four
defining-polynomial specifications. The resulting degree bounds are
$5059009872658432<2^{53}$ in 25 variables,
$99279020543180800<2^{57}$ in the fully specified 27-variable case,
and $43211777,39858177<2^{26}$ in the two 13-variable cases.
The explicit 27-variable coefficient height is less than $2^{2^{70}}$.
Variable counts, sufficient coefficient parameters, and arithmetic
premises are unchanged. The permitted retained-coordinate tuples are
enlarged, while the rational input projections remain the integers under
the corresponding stated premises.
\end{abstract}

\section{Results and arithmetic scope}
All variable counts include the distinguished rational input $X$.
Total degree and coefficient bounds concern finite specifications of
integer polynomials; the production determinants are not expanded.
Coefficient height is the maximum absolute integer coefficient after
the fixed coefficient parameter has been substituted.

\begin{theorem}\label{thm:main}
Under the arithmetic premises of the corresponding sources in
\cite{Unit,Subalgebra}, there are integer-polynomial specifications with
the following bounds:
\begin{center}\small
\begin{tabular}{lrrr}\toprule
Source & Variables & Degree at most & Strict bound\\\midrule
Inherited arithmetic premises &25&5059009872658432&$2^{53}$\\
Fully specified inherited-premise source &27&99279020543180800&$2^{57}$\\
Weak rational $abc$, exponent $121/100$ &13&43211777&$2^{26}$\\
Weak rational $abc$, exponent $16/15$ &13&39858177&$2^{26}$\\\bottomrule
\end{tabular}
\end{center}
Each polynomial has integer content one. The 25-variable family defines
$\Z$ for every sufficiently large fixed odd exponent $L$, with the same
sufficient threshold as its source. The 27-variable member at $L=20483$
defines $\Z$ and has coefficient height less than $2^{2^{70}}$.
The two 13-variable families define $\Z$ for the same sufficiently large
fixed odd shifts $B$ as their sources. The 25- and 13-variable sufficient
thresholds remain ineffective.
\end{theorem}

The general sources retain their elliptic group, height, normalized
modular realization, prime-pattern, and pole-realization premises.
The two conditional sources retain the indicated weak $abc$ inequality
and their specified remaining arithmetic inputs. No new realization
or height hypothesis is introduced. The full audit of these external
arithmetic premises remains outside the finite algebraic verification.
The smaller $L=3$ and $B=3$ computational members remain algebraic
fixtures, not numerically certified defining members.

\section{The pole obstruction with zero ordinates}
For a nonzero rational input $b$, put
\[
                         B=17b,\qquad V=B(B-b^2).
\]
The source pole equations are
\begin{equation}\label{eq:pole}
                 p^2=B(1-e)(1+3e),\qquad
                 Bq^2=(s^2-eV)(s^2+3eV).
\end{equation}
The zero-preparation formulation of~\cite[Section 3]{Dependencies}
allows zero witnesses. We give the relevant local argument explicitly.

\begin{lemma}\label{lem:pole}
If~\eqref{eq:pole} has a rational solution, with no nonzero condition
on $e,p,q,s$, then $v_\ell(b)$ is not a negative odd integer for any
prime $\ell\notin\{2,3,17\}$.
\end{lemma}
\begin{proof}
Suppose $\nu=v_\ell(b)<0$ is odd. Then $v_\ell(B)=\nu$ and
$v_\ell(V)=3\nu$, since the terms in $B-b^2$ have the distinct
valuations $\nu,2\nu$.
If $e=0$, the first right side is $B$, which cannot be a square.
If $v_\ell(e)>0$, its two displayed linear factors are units, again
giving an odd valuation. If $v_\ell(e)<0$, both factors have valuation
$v_\ell(e)$ and their product with $B$ has the odd valuation
$\nu+2v_\ell(e)$. Thus $v_\ell(e)=0$. This includes the possibilities
$e=1,-1/3$, when the first ordinate is zero.

Now $v_\ell(eV)=3\nu$ is odd. Its valuation differs from the even
valuation of $s^2$, with infinity allowed when $s=0$. Both factors in
the second right side are nonzero and have the same finite valuation
\[
                         \min\{2v_\ell(s),3\nu\}.
\]
Their product therefore has even valuation. The left side is zero if
$q=0$, and otherwise has odd valuation. Neither is possible.
\end{proof}

The rescaled pole equations in~\cite{Unit} have the same implication.
Their inverse coordinate changes divide only by the guarded input,
square multiplier, and conic denominator. They do not divide by the
second ordinate. Thus those maps remain valid when that ordinate is zero.

\section{Separation when the last radicand vanishes}
Let $A$ be a finite-dimensional commutative algebra over $\Q$ with a
quadratic tower presentation
\[
 A=\Q[t_1,\ldots,t_m],\qquad t_i^2=d_i,
 \qquad d_i\in\Q[t_1,\ldots,t_{i-1}],\qquad m\ge2.
\]
Here each quadratic adjunction is free of rank two; dependent square
classes are permitted. A geometric character is a unital homomorphism
$\chi:A\to\overline\Q$. Suppose every geometric value of $d_i$ is
nonzero for $i<m$. No such condition is imposed on $d_m$.
The final extension may therefore have nonzero nilpotents.

\begin{lemma}\label{lem:separation}
Choose positive rational numbers $M_i\ge1$ such that
$|\chi(t_i)|\le M_i$ for every geometric character, with a fixed
embedding into $\mathbb C$. Define backwards
\[
 \phi_{m+1}=0,\qquad
 \phi_i=\frac{t_i}{2M_i}+\frac{d_i}{4M_i^2}\phi_{i+1}.
\]
Then $|\chi(\phi_i)|\le2/3$ for every $i,\chi$.
The values $\chi(\phi_1)$ are nonzero and distinct as $\chi$ varies
over the geometric characters.
\end{lemma}
\begin{proof}
Backward induction gives the bound from
$1/2+(1/4)(2/3)=2/3$. At any character with $t_i\ne0$, the first
term in $\phi_i$ has magnitude $|t_i|/(2M_i)$, while the second has
magnitude at most $|t_i|^2/(6M_i^2)\le|t_i|/(6M_i)$.
In particular $\phi_1\ne0$, since $d_1\ne0$.

Let $\chi,\psi$ be distinct characters and let $i$ be the first
generator on which they differ. Their earlier values, and hence their
values of $d_i$, coincide. The two $t_i$ values are opposite and
nonzero; a zero value would give no difference. Their leading
contribution to $\phi_i(\chi)-\phi_i(\psi)$ has magnitude
$|t_i|/M_i$, while the remaining contribution has magnitude at most
$|d_i|/(3M_i^2)\le|t_i|/(3M_i)$. The difference is nonzero.
At each earlier stage it is multiplied by the same nonzero factor
$d_j/(4M_j^2)$, since $j<i\le m$. Thus the values of $\phi_1$
are distinct. The argument concerns characters, not their possible
multiplicities in a regular representation.
\end{proof}

Clearing denominators gives the unchanged integral separator formula
\begin{equation}\label{eq:omega}
 \Omega=\sum_{i=1}^m2M_it_i
             \left(\prod_{k<i}d_k\right)
             \left(\prod_{k>i}4M_k^2\right)
        =\left(\prod_{k=1}^m4M_k^2\right)\phi_1.
\end{equation}
The clearing factor is nonzero, so $\Omega$ has the same separation
and nonvanishing properties.

\begin{proposition}[Rational descent on a nonreduced algebra]
\label{prop:descent}
Let $A$ be any finite-dimensional commutative $\Q$-algebra. Suppose
$\Omega\in A$ has distinct nonzero values at its geometric characters,
and let $g\in\Q^\times$. If $M_\Omega$ denotes multiplication by
$\Omega$, then
\[
 \exists y\in\Q:\det(I-ygM_\Omega)=0
 \quad\Longleftrightarrow\quad
 \exists\chi:A\longrightarrow\Q\text{ a unital homomorphism}.
\]
\end{proposition}
\begin{proof}
Over $\overline\Q$, the algebra is a product of Artinian local
algebras. On a factor with residue character $\chi$, multiplication
by $\Omega$ is its scalar value $\chi(\Omega)$ plus a nilpotent
operator. Therefore
\[
 \det(I-ygM_\Omega)=
       \prod_{\chi}(1-yg\chi(\Omega))^{\ell_\chi},
\]
where $\ell_\chi\ge1$ is the dimension of that local factor.
If a rational $y$ makes the determinant zero, then $y\ne0$ and
some character has $\chi(\Omega)=1/(yg)\in\Q$.
Every Galois conjugate of $\chi$ has the same value on $\Omega$.
Separation makes it the same character. Thus every value of $\chi$
is fixed by $\operatorname{Gal}(\overline\Q/\Q)$ and belongs to $\Q$.
Conversely, a rational character gives $y=1/(g\chi(\Omega))$.
Nilpotents map to zero in every character and do not affect the argument.
No simple-eigenvalue hypothesis is used.
\end{proof}

Allowing an earlier radicand to vanish would not justify the same
conclusion for this formula. For example, in
$\Q[t,q]/(t^2,q^2-1)$, with $M_1=M_2=10$, formula~\eqref{eq:omega}
is $8000t$. It vanishes on both rational characters, and
$\det(I-YM_\Omega)=1$. This is why all preceding radicands remain units.

\section{The smaller scalar guard}
We retain the integral finite algebra, source circuits, quadratic tower,
root envelopes, and separator of~\cite{Subalgebra}. The point algebra
$\mathcal B$ is the tensor product of $p$ free rank-four point blocks;
the full algebra adjoins $s$ pole roots and has rank
\[
                            n=4^p2^s.
\]
The pairs $(p,s)$ are $(9,2),(9,3),(3,2),(3,2)$ in the four rows of
Theorem~\ref{thm:main}. Its rational point-algebra specializations are
always etale, as proved in~\cite{Unit}.

The source specifies integer circuits and guarded rational expressions.
Flattening multiplication and powers factors each circuit expression
into syntactic factors; no irreducible factorization is assumed.
Let $\mathcal F'$ contain the nonconstant factors of:
\begin{enumerate}
\item every native source guard except the final pole root;
\item every native pole radicand except the final one.
\end{enumerate}
Pole root symbols themselves are omitted from this list, since the
preceding radicands impose their nonvanishing. All source denominators
remain in the first group. Every factor of every preceding pole square
multiplier also remains in the second group.

For each $f\in\mathcal F'$, retain exactly its Galois-stable
multiquadratic subalgebra $\mathcal C_f$ specified in~\cite{Subalgebra}
and its integral norm $g_f$. Define
\begin{equation}\label{eq:guard}
                         G'=\prod_{f\in\mathcal F'}g_f.
\end{equation}
At every rational retained tuple, $G'\ne0$ means that all these factors
are units of the point algebra. In particular every denominator is
nonzero at every geometric character and all preceding pole radicands
are units. The last pole radicand is allowed to vanish.

For clarity, the terminal square extraction remains valid on this
larger locus. If the native equation is $Z^2=ND_0$ and
$ND_0=S^2c$, the reduced equation is $q_m^2=c$ and the recovered root
is $Z=S q_m$. Soundness uses multiplication by $S$, not division by it.
The native rational ordinate is $Z/D_0$, where $D_0$ remains guarded.
Neither $S$ nor $q_m$ needs to be nonzero. Existing full-unit
completeness witnesses have $S\ne0$ and are still represented by
$q_m=Z/S$.

The complete coefficient specification is
\begin{equation}\label{eq:polynomial}
                       P'=\det_{\mathcal A}(I-YG'M_\Omega).
\end{equation}
If $G'=0$, then $P'=1$, including at nonreduced specializations.
Otherwise Lemma~\ref{lem:separation} and
Proposition~\ref{prop:descent} give exactly the rational characters
of the enlarged guarded source. All the reconstructed pole equations
hold, with a possible zero final ordinate. Lemma~\ref{lem:pole}
gives the same local obstruction used throughout the source soundness
proofs. No other arithmetic formula, guard, or height estimate changes.
The same soundness arguments therefore prove integrality of the input.

Conversely each integer input has the old full-unit witnesses of
\cite{Unit}. The new guard is a subproduct of the old one, so these
witnesses remain allowed, and Proposition~\ref{prop:descent} supplies
the new rational $Y$. Thus completeness and its sufficient parameters
are unchanged. For the conditional sources use $XP'$ as before,
including input zero. Since $P'(\boldsymbol x,0)=1$, every final
polynomial has integer content one. This proves the semantic assertion
of Theorem~\ref{thm:main}. It is not an equivalence at every previously
retained base tuple: some zero-ordinate tuples are now permitted.

\section{Exact bounds and finite specifications}
Use the same weighted element degrees, coefficient norms, and integral
matrix basis as~\cite{Subalgebra}. Let $g,h$ bound respectively the
degree and binary logarithm of the coefficient sum of $G'$, and let
$o,l$ be the unchanged degree and logarithmic operator-norm bounds for
$\Omega$. The determinant estimate is
\begin{equation}\label{eq:bounds}
 \deg P'\le n(g+o+1),\qquad
 \log_2\|P'\|_1\le n(h+l+1),\qquad \deg_Y P'\le n.
\end{equation}
Multiplication by $X$ in the conditional cases adds one to total degree
and changes neither the coefficient norm nor the support cardinality.
These bounds do not require reduced specializations: they are estimates
for the same formal multiplication matrices over the integer coefficient
ring.

The exact majorants are
\begin{center}\scriptsize
\begin{tabular}{lrrrr}\toprule
Source & $g$ & $h$ & $o$ & $l$\\\midrule
25 &4824475361&33250472120&172420&1695856\\
27 &47339336484&392833791159924&589540&4941393118\\
13, $121/100$ &162135&520362&6660&29662\\
13, $16/15$ &150275&502572&5420&25322\\\bottomrule
\end{tabular}
\end{center}
The coefficient majorants for the 25-variable and both 13-variable rows
use the stored small fixtures. All degree majorants hold uniformly in the
allowed fixed parameter. The explicit 27-variable row uses $L=20483$ and gives
\[
 \log_2 H(P')\le823842533651079233536<2^{70},
\]
as asserted. No effective coefficient estimate for an unspecified
defining shift is inferred from a fixture.

For total variable count $V$ and $D_*=n(g+o+1)$, the support has the
upper bound
\[
 \binom{D_*+V}{V}-\binom{D_*-n+V-1}{V}.
\]
In the four rows these bounds are respectively less than
$2^{1193},2^{1401},2^{284},2^{283}$. These are upper bounds, not exact
expanded monomial counts.

The directory
\nolinkurl{computation/few_variables/terminal_root_v40}
contains four complete specifications, the generator, the independent
checker, and exact bound and verification reports. Each specification
pins its parent by SHA-256. The two removed native factor nodes are
\[
 \begin{array}{c|c}
 \text{source}&\text{removed nodes}\\\hline
 25&6223,6229\\
 27&7729,7735\\
 13\text{, either exponent}&1805,1811.
 \end{array}
\]
The checker reconstructs support from the native circuits, checks every
rational denominator, retains all preceding radicand factors, and
independently recomputes every final bound. It authenticates all
unchanged subalgebra and separator data through the separate parent
checker. Twenty deliberately altered specifications are rejected.
Eight small finite-algebra examples include negative and tiny radicands,
nonrational components, and terminal dual-number factors with repeated
rational eigenvalues. The displayed nonterminal-zero example checks the
restriction on which radicand may vanish. The mathematical descent and
arithmetic implications are proved above; the finite checks do not
certify the external arithmetic inputs.

\begin{thebibliography}{9}
\bibitem{Unit} \emph{A single norm determinant for the 25- and
13-variable constructions}, project manuscript,
\nolinkurl{papers/finite_algebra/unit_norm_20261114_v28.tex}.
\bibitem{Subalgebra} \emph{Smaller norm determinants for rational-point
encoding}, project manuscript,
\nolinkurl{papers/finite_algebra/subalgebra_norms_20261114_v30.tex}.
\bibitem{Dependencies} \emph{Arithmetic dependencies of the
defining-polynomial constructions}, project manuscript,
\nolinkurl{papers/arithmetic_dependencies/arithmetic_dependencies_20261108_v8.tex}.
The pole equations in its Section 3 expressly permit zero witnesses.
\end{thebibliography}
\end{document}
